\section{Balanceamento e ritmo de jogo}

No MVP, o jogo será um \textit{roguelike} de construção 
de mochilas, que poderá ser jogado repetidas vezes até 
que o jogador alcance o topo da torre. A dificuldade 
dos combates será ajustada por um modelo de 
\textit{Machine Learning}, com base no desempenho 
do jogador durante a escalada.

A maneira de combinar os itens deverá ser descoberta ao 
longo do jogo.

Na Fase 2, deverão ser implementadas mecânicas que 
limitem a obtenção rápida das chaves necessárias para 
acessar a Torre Pontiaguda e estabeleçam uma ordem 
específica para coletá-las, de maneira semelhante à 
progressão pelos ginásios da franquia Pokémon.

Além disso, uma quantidade adequada de combates deverá 
ser distribuída pelo mapa para fornecer dados que 
permitam treinar o modelo de \textit{Machine Learning}. 
O mapa será dividido em regiões de dificuldade, e a 
dificuldade inicial de cada região será definida com 
base nos resultados obtidos pelo jogador na região 
anterior.

Quando o jogador tiver acesso à Torre Pontiaguda, o 
desempenho obtido durante a progressão pelo mapa será 
utilizado como referência inicial para definir a 
dificuldade dos combates na torre.
